\chapter{Introduzione a Python}
\label{chap:introduzione-python}

Un algoritmo può essere tradotto in un programma mediante un linguaggio di
programmazione (\cref{sec:fi-algoritmi,sec:fi-linguaggi-esecuzione}). In questo
capitolo si considera Python 3, non pienamente compatibile con Python 2,
e se ne introducono le istruzioni fondamentali per gestire i dati,
controllare il flusso di esecuzione e organizzare il codice.\footnote{Materiale
di riferimento: S. Canazza, \emph{Python},
\nolinkurl{2026_FI_02_canazza_Python.pdf}.}

\section{Python e l'esecuzione dei programmi}
\label{sec:fi-python-esecuzione}

Python è un linguaggio \emph{ad alto livello} e \emph{general purpose}. Un
programma può essere salvato in un file sorgente, generalmente con estensione
\texttt{.py}, ed eseguito mediante un interprete. Nell'implementazione
CPython, il sorgente viene tradotto in \emph{bytecode}, poi eseguito
dalla macchina virtuale Python. La distinzione fra compilazione e
interpretazione è richiamata nella \cref{sec:fi-linguaggi-esecuzione}.

La funzione \texttt{print()} scrive sull'uscita standard. Una \emph{stringa}
è una sequenza di caratteri, delimitata nel codice da apici singoli o doppi;
il carattere \texttt{\#}, fuori dalle stringhe, introduce un commento fino
al termine della riga.

\begin{unipdcode}[language=Python]
# Un messaggio sul terminale
print("Ciao, mondo!")
\end{unipdcode}


Salvato il sorgente come \texttt{saluto.py}, lo si può eseguire dal
terminale:

\begin{unipdterminal}
$ python3 saluto.py
Ciao, mondo!
\end{unipdterminal}


Il nome del comando può essere \texttt{python} oppure
\texttt{python3}, a seconda dell'ambiente;
\verb|python3 --version| ne mostra la versione.
Gli esempi richiedono Python 3.9 o successivo.
Il simbolo \texttt{\$} nel terminale identifica il prompt e non fa parte
del comando.

Avviando l'interprete senza indicare un file si accede invece alla modalità
interattiva, riconoscibile dal prompt \lstinline|>>>|. Questa modalità
segue il ciclo \emph{REPL} (\emph{Read, Eval, Print, Loop}): valuta ogni
istruzione inserita e mostra il risultato delle espressioni, quando
previsto. In uno script, invece, le espressioni non vengono stampate
automaticamente.

\section{Variabili, tipi di dato e ingresso/uscita}
\label{sec:fi-python-variabili}

Per riutilizzare i valori durante l'esecuzione si impiegano le variabili.
L'\emph{assegnamento} associa un nome al valore di un'espressione:
\texttt{prezzo = 10} assegna \texttt{10} a \texttt{prezzo};
\texttt{prezzo = prezzo + 1} ne aggiorna il valore utilizzando quello
precedente. L'operatore \texttt{=} non indica quindi l'uguaglianza
matematica. Il riferimento a un nome non ancora definito provoca
\texttt{NameError}.

I nomi delle variabili, detti \emph{identificatori}, possono contenere
lettere, cifre e il carattere \texttt{\_}, ma non iniziare con una cifra.
Distinguono maiuscole e minuscole e non possono coincidere con parole
riservate, come \texttt{if} e \texttt{for}. Per i nomi composti si adotta
generalmente la convenzione \emph{snake\_case}, ad esempio
\texttt{prezzo\_unitario}.

In Python il tipo appartiene al valore, non al nome: un nome può essere
riassegnato a valori di tipi diversi. I tipi fondamentali sono
\texttt{int} per gli interi, \texttt{float} per i numeri in virgola mobile,
\texttt{str} per le stringhe e \texttt{bool} per i valori
\texttt{True} e \texttt{False}. La funzione \texttt{type()} restituisce
il tipo di un valore. Gli \texttt{int} non hanno una larghezza fissa,
mentre le approssimazioni dei \texttt{float} dipendono dalla
rappresentazione descritta nella
\cref{sec:fi-precisione-arrotondamento}.

Gli operatori \texttt{+}, \texttt{-}, \texttt{*} e \texttt{**} indicano
rispettivamente somma, sottrazione, prodotto e potenza. Per la divisione
si distinguono:
\[
\begin{array}{rcl@{\qquad}l}
  \texttt{17 / 5}  &=& 3.4 & \text{divisione ordinaria},\\
  \texttt{17 // 5} &=& 3 & \text{quoziente arrotondato verso }-\infty,\\
  \texttt{17 \% 5}  &=& 2 & \text{resto della divisione}.
\end{array}
\]
La divisione \texttt{/} produce un valore \texttt{float} anche quando il
quoziente è intero; \texttt{//} arrotonda verso il basso e non sempre
verso zero. La potenza precede la moltiplicazione, che precede
l'addizione; le parentesi permettono di stabilire esplicitamente l'ordine
delle operazioni.

Il significato degli operatori dipende dai tipi degli operandi: per le
stringhe, \texttt{+} esegue la concatenazione, mentre la moltiplicazione
per un intero ne ripete il contenuto. Ad esempio,
\lstinline|"3" + "4"| produce \lstinline|"34"|, non \texttt{7};
tentare di sommare direttamente una stringa e un intero provoca
\texttt{TypeError}. Le funzioni \texttt{int()}, \texttt{float()},
\texttt{str()} e \texttt{bool()} eseguono conversioni esplicite.
In particolare, \texttt{int(3.9)} restituisce \texttt{3}, poiché
la conversione tronca verso zero. Inoltre, \texttt{bool(0)} restituisce
\texttt{False} e \texttt{bool(3)} restituisce \texttt{True}.
Una conversione non ammessa, come \lstinline|int("ciao")|,
genera \texttt{ValueError}.

La funzione \texttt{input()} legge una riga dall'ingresso standard e
restituisce una stringa, anche se l'utente ha digitato un numero.
Per eseguire calcoli occorre quindi convertire i valori acquisiti.

\begin{unipdcode}[language=Python]
primo = int(input("Primo intero: "))
secondo = int(input("Secondo intero: "))
totale = primo
totale = totale + secondo
print("Somma:", totale)
print("Media:", totale / 2)
\end{unipdcode}


Il programma assume che le due stringhe inserite rappresentino
numeri interi validi.

\section{Logica booleana e istruzioni condizionali}
\label{sec:fi-python-condizioni}

Per esprimere le condizioni che regolano il flusso di esecuzione si
utilizzano espressioni booleane. Gli operatori di confronto
\texttt{==}, \texttt{!=}, \texttt{<}, \texttt{<=}, \texttt{>} e
\texttt{>=} restituiscono valori booleani.
In particolare, \texttt{==} verifica l'uguaglianza, mentre \texttt{=}
esegue un assegnamento.

Gli operatori logici \texttt{and}, \texttt{or} e \texttt{not} esprimono
congiunzione, disgiunzione e negazione. Nelle seguenti tabelle di verità,
\(\mathrm{V}\) e \(\mathrm{F}\) indicano rispettivamente
\texttt{True} e \texttt{False}:
\[
\begin{array}{cc|cc}
  a & b & a\mathbin{\texttt{ and }}b & a\mathbin{\texttt{ or }}b\\
  \hline
  \mathrm{F} & \mathrm{F} & \mathrm{F} & \mathrm{F}\\
  \mathrm{F} & \mathrm{V} & \mathrm{F} & \mathrm{V}\\
  \mathrm{V} & \mathrm{F} & \mathrm{F} & \mathrm{V}\\
  \mathrm{V} & \mathrm{V} & \mathrm{V} & \mathrm{V}
\end{array}
\qquad
\begin{array}{c|c}
  a & \texttt{not }a\\
  \hline
  \mathrm{F} & \mathrm{V}\\
  \mathrm{V} & \mathrm{F}
\end{array}
\]
La precedenza è \texttt{not}, poi \texttt{and}, infine \texttt{or}.
Le \emph{leggi di De Morgan} forniscono le equivalenze
\[
\begin{aligned}
 \texttt{not (a and b)}&\quad\Longleftrightarrow\quad
 \texttt{(not a) or (not b)},\\
 \texttt{not (a or b)}&\quad\Longleftrightarrow\quad
 \texttt{(not a) and (not b)}.
\end{aligned}
\]
Gli operatori \texttt{and} e \texttt{or} valutano il secondo operando
soltanto quando necessario (\emph{valutazione a cortocircuito}).

L'istruzione \texttt{if} esegue un blocco quando la condizione è vera;
\texttt{elif} introduce condizioni successive e \texttt{else} il caso
residuo. In una catena condizionale viene eseguito soltanto il blocco
associato alla prima condizione soddisfatta.

\begin{unipdcode}[language=Python]
voto = 27
if voto < 18:
    print("Insufficiente")
elif voto < 24:
    print("Sufficiente")
elif voto < 28:
    print("Buono")
else:
    print("Ottimo")
\end{unipdcode}


L'\emph{indentazione} identifica i blocchi di istruzioni ed è parte
della sintassi di Python. Dopo i due punti, le istruzioni del blocco
sono rientrate allo stesso livello; terminato il blocco, si ritorna
all'indentazione precedente. La convenzione è utilizzare quattro spazi
per livello. Un rientro non valido può produrre \texttt{IndentationError};
un rientro sintatticamente valido ma errato può invece modificare la
logica del programma.

\section{Cicli e iterazioni}
\label{sec:fi-python-cicli}

Le istruzioni condizionali selezionano il blocco da eseguire; un
\emph{ciclo} permette invece di ripeterlo. Ogni esecuzione del blocco
è detta \emph{iterazione}. Il ciclo \texttt{while} controlla una condizione
prima di ogni iterazione e termina quando essa diventa falsa. Il corpo
può quindi non essere eseguito neppure una volta.

\begin{unipdcode}[language=Python]
eta = int(input("Eta non negativa: "))
while eta < 0:
    eta = int(input("Valore non valido, riprovare: "))
print("Eta accettata:", eta)
\end{unipdcode}


L'esempio richiede nuovamente l'età finché il valore è negativo. La
terminazione dipende dall'inserimento di un intero non negativo; una
stringa non convertibile provocherebbe \texttt{ValueError}. Un ciclo \texttt{while} può non terminare se la condizione
rimane sempre vera.

Nei \emph{cicli con contatore} si distinguono inizializzazione,
condizione e aggiornamento. L'istruzione \texttt{n += 1} abbrevia
\texttt{n = n + 1}. Nei \emph{cicli con sentinella}, invece, un valore
prestabilito segnala la fine dell'acquisizione e non viene elaborato come
dato.

\begin{unipdcode}[language=Python]
somma = 0
numero = int(input("Numero (0 per terminare): "))
while numero != 0:
    somma += numero
    numero = int(input("Numero (0 per terminare): "))
print("Somma:", somma)
\end{unipdcode}


L'inserimento di zero termina il ciclo senza contribuire alla somma;
se è il primo valore, la somma resta nulla.

Il ciclo \texttt{for} visita, nell'ordine, i valori di un oggetto
\emph{iterabile}, assegnandoli alla variabile indicata. La funzione
\texttt{range()} fornisce una sequenza di interi con estremo iniziale
incluso ed estremo finale escluso:
\[
\begin{array}{ll}
  \texttt{range(5)} & 0,1,2,3,4,\\
  \texttt{range(2,6)} & 2,3,4,5,\\
  \texttt{range(0,10,2)} & 0,2,4,6,8,\\
  \texttt{range(5,0,-1)} & 5,4,3,2,1.
\end{array}
\]
Con tre argomenti \texttt{range(inizio, fine, passo)}, il passo può
essere negativo, ma non nullo. Il ciclo \texttt{for} non richiede
l'aggiornamento manuale della variabile di iterazione.

\begin{unipdcode}[language=Python]
totale = 0
for numero in range(1, 6):
    totale = totale + numero
print("Somma:", totale)
\end{unipdcode}


La variabile \texttt{totale} è un \emph{accumulatore}: contiene la somma
parziale e viene aggiornata a ogni iterazione. In questo caso il ciclo
somma gli interi da uno a cinque, individuati da \texttt{range(1, 6)},
ottenendo \texttt{15}. La ricorsione fornisce un'altra forma di
ripetizione del calcolo, trattata nella
\cref{sec:fi-ricorsione-terminazione}.

\section{Liste}
\label{sec:fi-python-liste}

I cicli possono elaborare anche sequenze di dati. Una \emph{lista} è
una sequenza modificabile di elementi, scritta nel codice tra
parentesi quadre. Gli elementi possono avere tipi differenti e sono
individuati dalla loro posizione, a partire dall'indice zero. \texttt{len(lista)} restituisce il numero di elementi:
per una lista di lunghezza \(n\), gli indici non negativi validi sono
\(0,\ldots,n-1\); \texttt{lista[-1]} indica l'ultimo elemento.
Un indice non valido provoca \texttt{IndexError}.

L'assegnamento \texttt{lista[i] = valore} modifica un elemento.
Il metodo \texttt{append(valore)} aggiunge un elemento in fondo;
\texttt{pop()} rimuove e restituisce l'ultimo. L'operatore \texttt{in}
verifica la presenza di un valore.

Lo \emph{slicing} \texttt{lista[i:j]} estrae una nuova lista con estremo
iniziale incluso e finale escluso. \texttt{sorted(lista)} restituisce
una nuova lista ordinata senza modificare quella originale, se gli
elementi sono confrontabili. Gli algoritmi impiegati per ordinare
i dati sono studiati nel \cref{chap:ricerca-ordinamento}.

\begin{unipdcode}[language=Python]
voti = [20, 30, 27]
voti[0] = 21
voti.append(18)
ultimo = voti.pop()
ordinati = sorted(voti)
print("Primi due:", voti[:2])
print("Presente 30:", 30 in voti)

totale = 0
conteggio = 0
for voto in voti:
    if voto >= 24:
        totale = totale + voto
        conteggio = conteggio + 1
if conteggio > 0:
    print("Media dei voti almeno 24:", totale / conteggio)
\end{unipdcode}


L'esempio combina accessi agli elementi, modifica della lista, ricerca,
ordinamento e accumulazione. Il ciclo visita i voti e calcola la media
di quelli non inferiori a \texttt{24}; il controllo
\texttt{if conteggio > 0} evita la divisione per zero quando nessun
valore soddisfa la condizione.

\section{Funzioni}
\label{sec:fi-python-funzioni}

Per organizzare operazioni riutilizzabili si impiegano le \emph{funzioni},
che associano un nome a un blocco di istruzioni. Una funzione si definisce
mediante \texttt{def}; la definizione non esegue il corpo, che viene
eseguito soltanto quando la funzione è \emph{chiamata}. I nomi dichiarati nella definizione sono i
\emph{parametri}; i valori forniti alla chiamata sono gli
\emph{argomenti}.

\begin{unipdcode}[language=Python]
def area_rettangolo(base, altezza):
    return base * altezza


risultato = area_rettangolo(8, 5)
print("Area:", risultato)
\end{unipdcode}


L'istruzione \texttt{return} termina la chiamata e restituisce il
valore calcolato al chiamante. È distinta da \texttt{print()}, che
scrive sull'uscita standard: una funzione che termina senza
\texttt{return} restituisce \texttt{None}. I parametri e le variabili
assegnate nel corpo della funzione sono di norma \emph{locali} alla
chiamata. Una funzione può restituire un valore booleano, come
nell'espressione \texttt{return numero \% 2 == 0}.

Le definizioni possono essere raccolte in \emph{moduli}, ossia file
Python importabili. Salvando la definizione di
\texttt{area\_rettangolo} in \texttt{geometria.py}, un secondo file
può utilizzarla così:

\begin{unipdcode}[language=Python]
import geometria

print("Area:", geometria.area_rettangolo(8, 5))
\end{unipdcode}


L'istruzione \texttt{import geometria} rende disponibile il modulo;
la chiamata \texttt{geometria.area\_rettangolo(8, 5)} ne usa la funzione.
L'estensione \texttt{.py} non compare in \texttt{import}. Le funzioni
possono anche richiamare se stesse, come si vedrà nel
\cref{chap:ricorsione}.

\section{Gestione dei file di testo}
\label{sec:fi-python-file}

Oltre al codice organizzato in moduli, un programma può utilizzare file
per conservare dati tra esecuzioni successive. I file sono memorizzati
nella memoria di massa (\cref{sec:fi-memoria-archiviazione}). In un
\emph{file di testo}, i caratteri sono memorizzati secondo una codifica,
come UTF-8 (\cref{sec:fi-caratteri}). Quando il nome non contiene un percorso
assoluto, il file viene cercato nella directory di lavoro corrente,
che non coincide necessariamente con quella del programma.

La funzione \texttt{open()} apre il file. Le modalità principali sono
\lstinline|"r"| per leggere, \lstinline|"w"| per creare o sovrascrivere
e \lstinline|"a"| per aggiungere dati in fondo. La modalità
\lstinline|"w"| cancella il contenuto preesistente all'apertura;
\lstinline|"r"| richiede invece che il file esista.

Il costrutto \texttt{with open(...) as file:} chiude il file
all'uscita dal blocco, anche se si verifica un errore. È opportuno
specificare la codifica UTF-8 mediante il parametro \texttt{encoding}.

\clearpage

\begin{unipdcode}[language=Python]
with open("voti.txt", "w", encoding="utf-8") as file:
    file.write("24\n27\n")

with open("voti.txt", "a", encoding="utf-8") as file:
    file.write("30\n")
\end{unipdcode}


Il metodo \texttt{write()} scrive una stringa senza aggiungere
automaticamente un fine riga; la sequenza \lstinline|"\n"| rappresenta
un carattere di nuova riga. L'esempio produce il file \texttt{voti.txt}
contenente tre interi, uno per riga. La prima apertura in modalità
\lstinline|"w"| sovrascrive eventuali dati precedenti; la seconda,
in modalità \lstinline|"a"|, aggiunge l'ultima riga.

Il metodo \texttt{read()} restituisce l'intero contenuto del file
come stringa. In alternativa, un ciclo \texttt{for} permette di leggerlo
\emph{riga per riga}; ogni riga contiene normalmente anche il
carattere di fine riga. \texttt{strip()} rimuove spazi bianchi iniziali
e finali, compresi i terminatori di riga. Un numero letto da un file
rimane testo finché non viene convertito con \texttt{int()} o
\texttt{float()}.

\begin{unipdcode}[language=Python]
try:
    totale = 0
    with open("voti.txt", "r", encoding="utf-8") as file:
        for riga in file:
            totale = totale + int(riga.strip())
    print("Somma:", totale)
except FileNotFoundError:
    print("File voti.txt non trovato.")
\end{unipdcode}


Il programma stampa \texttt{Somma: 81} se il file contiene i tre
valori dell'esempio precedente. La clausola
\texttt{except FileNotFoundError} gestisce l'assenza del file quando
questo viene aperto in lettura; contenuti non convertibili in interi
provocherebbero invece \texttt{ValueError}.
